% ----------------------------------------------------------------
% Section: Results and Evaluation
% ----------------------------------------------------------------
\chapter{RESULTS AND EVALUATION}
\label{ch:results}

\section{Overview}

This chapter presents the results obtained from executing the complete build pipeline and booting the resulting image on a Raspberry Pi~5 Model B (16~GB RAM variant). All observations were made via UART serial console at 115200~baud using \texttt{minicom}. The serial interface connects to GPIO~14 (TX) and GPIO~15 (RX) on the 40-pin header \cite{rpi_uart_pinout}.

% ----------------------------------------------------------------
\section{Build Pipeline Execution}
% ----------------------------------------------------------------

The full pipeline was executed using:

\begin{lstlisting}[language=bash]
rpi5-build --config config/rpi5_uboot_aosp_signed.yaml \
           --stage all --code all
\end{lstlisting}

\subsection{U-Boot Build}

U-Boot was compiled from commit \texttt{eefb822fb574} (pinned in the config) using \path{rpi_arm64_defconfig}. After AVB Kconfig injection, the \texttt{.config} contained all required symbols:

\begin{lstlisting}
CONFIG_ANDROID_BOOT_IMAGE=y
CONFIG_FASTBOOT=y
CONFIG_LIBAVB=y
CONFIG_AVB_VERIFY=y
CONFIG_CMD_AVB=y
CONFIG_FASTBOOT_BUF_ADDR=0x10000000
CONFIG_AVB_BUF_ADDR=0x10000000
\end{lstlisting}

The AVB root public key (1032~bytes, SHA-256/RSA-4096 packed format) was injected into \texttt{common/avb\_verify.c} and the build completed successfully.

\subsection{AOSP Build}

The AOSP tree was built with \texttt{RPI5\_ENABLE\_AVB=true RPI5\_ENABLE\_FBE=true}, producing:

\begin{itemize}
    \item \texttt{boot.img} --- Android kernel + ramdisk with \texttt{fstab.rpi5.fbe}
    \item \texttt{system.img} --- ext4 system partition (pre-signing)
    \item \texttt{vendor.img} --- ext4 vendor partition (pre-signing)
\end{itemize}

\subsection{Signing Stage}

The sign stage processed both partitions:

\begin{itemize}
    \item \texttt{system.img}: resized to fit the 3072~MB AVB partition limit, hashtree footer appended.
    \item \texttt{vendor.img}: resized to fit the 384~MB AVB partition limit, hashtree footer appended.
    \item \texttt{vbmeta.img}: generated referencing both signed images.
\end{itemize}

\subsection{SD Card Image Assembly}

The signed image (\texttt{rpi5-uboot-aosp-signed-user-202606040454.img}, 14.3~GiB) was flashed to a 64~GiB SD card using \texttt{scripts/flash\_sdcard.sh}. Because the image is built for a fixed partition layout smaller than the target card, the script performs a three-step post-flash expansion:

\begin{enumerate}
    \item \textbf{Relocate backup GPT header} to the physical end of the 64~GiB disk.
    \item \textbf{Expand \texttt{userdata} (p3)} to fill the remaining space, growing from its reserved size to 56.0~GiB.
    \item \textbf{Resize the ext4 filesystem} on \texttt{/dev/sdc3} to match the new partition boundary.
\end{enumerate}

The resulting partition layout confirms all six partitions from Table~\ref{tab:partitions} are correctly placed and sized:

\begin{lstlisting}[caption={Final GPT partition layout after flashing to 64~GiB SD card}]
Number  Start (sector)  End (sector)   Size        Name
   1           2048        264191    128.0 MiB    boot
   3        7383040     124735454     56.0 GiB    userdata  (expanded)
   5         264192       6555647      3.0 GiB    system
   6        6555648       7342079    384.0 MiB    vendor
   7        7342080       7374847     16.0 MiB    metadata
   8        7374848       7383039      4.0 MiB    vbmeta
\end{lstlisting}

The \texttt{userdata} partition fills whatever space remains on the card, making the pipeline card-size agnostic: the same signed image can be flashed to any SD card of 14.3~GiB or larger without rebuilding.

% ----------------------------------------------------------------
\section{Boot Sequence Validation}
% ----------------------------------------------------------------

Boot validation was performed across two image variants: an unsigned \texttt{eng} build (no AVB, SELinux permissive, \texttt{fail\_open} policy) and a signed \texttt{user} build (AVB enabled, SELinux enforcing, \texttt{fail\_closed} policy). The unsigned build provides a functional baseline; the signed build validates the security stack. Negative test cases within each group confirm that the security mechanisms reject invalid inputs rather than silently accepting them.

% ----------------------------------------------------------------
\subsection{Unsigned Build Validation}
% ----------------------------------------------------------------

The unsigned build uses a plain U-Boot boot script without AVB commands, an MBR partition layout (no vbmeta partition), and \texttt{androidboot.selinux=permissive}. Five test cases validate its expected behaviour.

\paragraph{TC-U-001 --- Unsigned image boots to Android.}
U-Boot loaded the kernel and ramdisk from the FAT32 boot partition and transferred control to the kernel without invoking any AVB commands:

\begin{lstlisting}[caption={TC-U-001: U-Boot output on unsigned image}]
** Booting bootflow 'mmc@fff000.bootdev.part_1' with script
=== Android U-Boot Boot Script ===
Loading kernel Image...
37157376 bytes read in 1577 ms (22.5 MiB/s)
Loading ramdisk...
1544522 bytes read in 68 ms (21.7 MiB/s)
Booting Android...
\end{lstlisting}

Android completed initialisation and the serial console became responsive, confirming a full unsigned boot.

\paragraph{TC-U-002 --- Boot partition contains correct files for unsigned image.}
The FAT32 boot partition (p1) was inspected on the host before flashing. The plain (non-AVB) boot script was present and \texttt{boot\_avb.scr} was absent, confirming the unsigned pipeline produced the correct boot artefacts:

\begin{lstlisting}[caption={TC-U-002: Boot partition file listing --- unsigned image}]
-rwxr-xr-x  boot.scr       1,499  (non-AVB variant)
-rwxr-xr-x  cmdline.txt      124
-rwxr-xr-x  config.txt       642
-rwxr-xr-x  Image         37,157,376
-rwxr-xr-x  ramdisk.img    1,544,522
-rwxr-xr-x  u-boot.bin       708,720
\end{lstlisting}

The \texttt{boot.scr} size (1\,499~bytes) and \texttt{u-boot.bin} size (708\,720~bytes) are characteristic of the non-AVB variants; both are smaller than their signed counterparts (see TC-S-002), as neither the AVB boot script nor \texttt{libavb} are compiled in.

\paragraph{TC-U-003 --- AVB verification absent from unsigned boot path.}
The U-Boot serial output in TC-U-001 contains no occurrences of \texttt{avb init}, \texttt{avb verify}, \texttt{Initializing AVB}, \texttt{Verifying partitions}, or \texttt{AVB:}.  Execution proceeds directly from \texttt{Loading kernel Image\ldots} to \texttt{Booting Android\ldots} to \texttt{booti}, confirming that the unsigned image never invokes the AVB verification path and provides no cryptographic boot-integrity check.

\paragraph{TC-U-004 --- First-stage init reports absent vbmeta partition (negative).}
The unsigned image uses an MBR partition layout with no p8 (vbmeta) partition.  Android's first-stage init attempted to locate vbmeta and reported its absence:

\begin{lstlisting}[caption={TC-U-004: First-stage init reporting absent vbmeta partition}]
[    4.126439] init: Failed to read vbmeta partitions.
\end{lstlisting}

This confirms that first-stage init correctly identifies the absence of vbmeta and that the unsigned image carries no AVB integrity guarantees.

\paragraph{TC-U-005 --- Boot diagnostics service reports unsigned security state.}
The \texttt{rpi5\_boot\_diag} service fires on \texttt{sys.boot\_completed=1} and prints a structured security summary to the UART console.  On the unsigned build it produced:

\begin{lstlisting}[caption={TC-U-005: rpi5\_boot\_diag.sh output --- unsigned build}]
 rpi5 boot diagnostics  Sun May  3 06:08:12 GMT 2026
============================================================

--- verified boot ---
orange
unlocked
1.1

--- global crypto state ---
unsupported

--- selinux ---
Permissive
\end{lstlisting}

The state \texttt{orange}/\texttt{unlocked} is the expected verified boot state for a build without a locked bootloader.  \texttt{crypto state: unsupported} confirms no FDE/FBE is active.  \texttt{SELinux: Permissive} matches the unsigned build configuration.

% ---- Unsigned summary table ----
\begin{table}[h]
\centering
\caption{Unsigned build validation results}
\label{tab:tc_unsigned}
\begin{tabular}{p{2.2cm}p{2.0cm}p{7.8cm}p{1.4cm}}
\toprule
\textbf{TC ID} & \textbf{Type} & \textbf{Description} & \textbf{Result} \\
\midrule
TC-U-001 & Functional & Unsigned image boots to Android without AVB & \checkmark Pass \\
TC-U-002 & Functional & Boot partition contains correct files (non-AVB variants) & \checkmark Pass \\
TC-U-003 & Security   & No AVB commands present on unsigned boot path & \checkmark Pass \\
TC-U-004 & Security   & First-stage init reports absent vbmeta partition & \checkmark Pass \\
TC-U-005 & Functional & Boot diagnostics confirms unsigned security state & \checkmark Pass \\
\bottomrule
\end{tabular}
\end{table}

% ----------------------------------------------------------------
\subsection{Signed Build Validation}
% ----------------------------------------------------------------

The signed build uses an AVB-enabled U-Boot, a GPT partition layout with a dedicated vbmeta partition (p8), \texttt{fail\_closed} policy, and \texttt{androidboot.selinux=enforcing}.  Seven test cases validate its security behaviour, including two negative cases that confirm fail-closed rejection.

\paragraph{TC-S-001 --- Signed image boots with AVB verification passed.}
With the PL011-AXI patch applied, U-Boot produced serial output on power-on.  The AVB boot script ran, verified the vbmeta signature, and transferred control to Android:

\begin{lstlisting}[caption={TC-S-001: U-Boot AVB boot script output --- signed build (abbreviated)}]
=== Android U-Boot Boot Script (AVB) ===
Setting AVB policy to default: fail_closed
AVB fail policy: fail_closed
Initializing AVB...
Verifying partitions...
## Android Verified Boot 2.0 version 1.1.0
Verification passed successfully
AVB: Verification PASSED
Loading kernel Image...
Loading ramdisk...
bootargs: root=/dev/ram0 rootwait androidboot.hardware=rpi5
  androidboot.selinux=permissive androidboot.encryption_mode=fbe
  androidboot.vbmeta.device=PARTUUID=66680148-de99-4703-8f30-76949b3cec9a
  androidboot.vbmeta.avb_version=1.1
  androidboot.vbmeta.device_state=unlocked
  androidboot.vbmeta.hash_alg=sha256
  androidboot.vbmeta.digest=68736b3481643cfadb7d69022e6ae6a59370236695b3fe64c14daf87ee36f86e
  androidboot.veritymode=enforcing
  androidboot.verifiedbootstate=orange
Booting Android with AVB verification...
\end{lstlisting}

\texttt{androidboot.verifiedbootstate=orange} indicates that AVB signature verification passed but the bootloader is in unlocked state --- the expected outcome on RPi5, which has no OEM fuse-locked key store \cite{avb2_readme}.  \texttt{androidboot.veritymode=enforcing} confirms that dm-verity will enforce partition integrity at runtime.

\paragraph{TC-S-002 --- Boot partition contains correct files for signed image.}
The FAT32 boot partition (p1) of the signed image was inspected on the host.  The AVB boot script was present and both \texttt{boot.scr} and \texttt{u-boot.bin} were larger than their unsigned counterparts, reflecting the AVB commands and embedded \texttt{libavb} respectively:

\begin{lstlisting}[caption={TC-S-002: Boot partition file listing --- signed image}]
-rwxr-xr-x  boot.scr       3,289  (AVB variant -- larger than unsigned 1,499)
-rwxr-xr-x  cmdline.txt      176
-rwxr-xr-x  config.txt       642
-rwxr-xr-x  Image         37,157,376
-rwxr-xr-x  ramdisk.img    1,326,049
-rwxr-xr-x  u-boot.bin       769,928  (AVB-enabled -- larger than unsigned 708,720)
\end{lstlisting}

The size increase in \texttt{u-boot.bin} (769\,928 vs 708\,720~bytes) reflects the inclusion of \texttt{libavb} and the embedded RSA-4096 public key.  The size increase in \texttt{boot.scr} (3\,289 vs 1\,499~bytes) reflects the AVB \texttt{avb init}/\texttt{avb verify} command sequence.

\paragraph{TC-S-003 --- Kernel boot-layer markers confirm full boot chain.}
Four custom diagnostic markers were visible in the serial log, confirming each stage of the boot chain from kernel entry to userspace handoff:

\begin{lstlisting}[caption={TC-S-003: Kernel RPi5-DIAG boot-layer markers (abbreviated)}]
[    0.000000][    T0] RPi5-DIAG: Layer 1 start_kernel entry (bootloader handoff complete)
[    0.000000][    T0] RPi5-DIAG: Layer 2 console_init complete (early console up, mm init done)
[    0.009196][    T1] RPi5-DIAG: Layer 3 do_initcalls entry (driver probe sequence begins)
[   22.617616][    T1] RPi5-DIAG: Layer 4 kernel_init userspace exec (all drivers probed, exec /init)
\end{lstlisting}

Layer~4 appears at approximately 22.6~seconds from kernel entry, reflecting the time for all driver probes and module initialisations on the BCM2712.

\paragraph{TC-S-004 --- dm-verity active on system and vendor partitions.}
First-stage init mounted \texttt{/system} and \texttt{/vendor} through device-mapper rather than directly from the raw block devices, confirming dm-verity is active on both partitions:

\begin{lstlisting}[caption={TC-S-004: First-stage init dm-verity mount messages (abbreviated)}]
[   22.698811][    T1] init: [libfs_mgr] check_fs():
    mount(/dev/block/by-name/metadata,/metadata,ext4)=0: Success
[   22.815684][    T1] EXT4-fs (dm-0): mounted filesystem
    d392d8bb-6c67-58f7-b0b2-091ad46c5bc2 ro without journal. Quota mode: none.
[   22.854598][    T1] EXT4-fs (dm-1): mounted filesystem
    2b96c597-1e2f-5ee1-9851-c4a9fa9de36e ro without journal. Quota mode: none.
\end{lstlisting}

The \texttt{dm-0} and \texttt{dm-1} device names confirm that \texttt{/system} and \texttt{/vendor} are served through device-mapper rather than directly from \texttt{mmcblk0p5}/\texttt{p6}.

\paragraph{TC-S-005 --- FBE initialisation on signed build (partial).}
On the signed build with FBE enabled (\texttt{RPI5\_ENABLE\_FBE=true}), AVB verification passed but \texttt{vold} failed to initialise the FBE key directory before \texttt{sys.boot\_completed=1} was reached, causing a reboot loop at approximately 9~seconds:

\begin{lstlisting}[caption={TC-S-005: FBE boot loop on signed build}]
AVB: Verification PASSED
...
[    2.893300][  T350] vold: Failed to create /metadata/vold: Permission denied
[    2.900410][  T350] vold: read_key failed in mountFstab
[    9.306520][    T1] reboot: Restarting system with command 'bootloader'
\end{lstlisting}

The root cause (unlabelled \texttt{metadata} partition root inode under SELinux Enforcing mode) and its fix (\texttt{e2fsdroid} replacing bare \texttt{mkfs.ext4} in the SD card stage) are documented in Section~\ref{sec:impl_fbe_selinux}.  End-to-end FBE validation could not be completed within the project timeline.

\paragraph{TC-S-006 --- Wrong public key rejected by fail-closed policy (negative).}
A U-Boot binary compiled with a different AVB root public key was placed on a signed SD card, with the signed vbmeta and partition images left unchanged.  This scenario simulates a firmware-substitution attack.  \texttt{libavb} detected the key mismatch and U-Boot reported:

\begin{lstlisting}[caption={TC-S-006: U-Boot wrong-key rejection}]
avb_slot_verify.c:837: ERROR: vbmeta: Public key used to sign data rejected.
Verification failed, reason: Public keys are not accepted
AVB: Verification FAILED
Fail-closed policy: resetting
resetting ...
\end{lstlisting}

The device reset immediately without loading the kernel, confirming that the fail-closed policy prevents boot when the trust anchor does not match.

\paragraph{TC-S-007 --- Corrupt vbmeta header rejected by fail-closed policy (negative).}
The first 256~bytes of the vbmeta partition (p8) were zeroed, destroying the \texttt{AVB0} magic.  The rejection occurred at the header-parse stage, earlier in the verification flow than the key-check rejection in TC-S-006.  U-Boot reported:

\begin{lstlisting}[caption={TC-S-007: U-Boot vbmeta header rejection}]
avb_vbmeta_image.c:46: ERROR: Magic is incorrect.
avb_slot_verify.c:745: ERROR: vbmeta: Error verifying vbmeta image: invalid vbmeta header
Verification failed, reason: Metadata is invalid or inconsistent
AVB: Verification FAILED
Fail-closed policy: resetting
resetting ...
\end{lstlisting}

The device reset immediately without loading the kernel, confirming that vbmeta integrity is enforced before any signature or key check is reached.

% ---- Signed summary table ----
\begin{table}[h]
\centering
\caption{Signed build validation results}
\label{tab:tc_signed}
\begin{tabular}{p{2.2cm}p{2.0cm}p{7.4cm}p{1.8cm}}
\toprule
\textbf{TC ID} & \textbf{Type} & \textbf{Description} & \textbf{Result} \\
\midrule
TC-S-001 & Functional & Signed image boots with AVB verification passed & \checkmark Pass \\
TC-S-002 & Functional & Boot partition contains correct files (AVB variants) & \checkmark Pass \\
TC-S-003 & Functional & Kernel boot-layer markers confirm full boot chain & \checkmark Pass \\
TC-S-004 & Security   & dm-verity active on system and vendor partitions & \checkmark Pass \\
TC-S-005 & Security   & FBE initialisation on signed build & $\sim$ Partial \\
TC-S-006 & Security   & Wrong public key rejected by fail-closed policy & \checkmark Pass \\
TC-S-007 & Security   & Corrupt vbmeta header rejected by fail-closed policy & \checkmark Pass \\
\bottomrule
\end{tabular}
\end{table}

% ----------------------------------------------------------------
\section{Evaluation Summary}
% ----------------------------------------------------------------

Table~\ref{tab:eval} summarises the evaluation against the objectives defined in Chapter~\ref{ch:introduction}.

\begin{table}[h]
\centering
\caption{Evaluation against project objectives}
\label{tab:eval}
\begin{tabular}{lll}
\toprule
\textbf{Objective} & \textbf{Status} & \textbf{Evidence} \\
\midrule
AVB~2.0 evaluation  & \checkmark Pass    & \texttt{verifiedbootstate=orange}; \texttt{veritymode=enforcing}; \\
                    &                    & \quad fail-closed policy confirmed (TC-S-006, TC-S-007) \\
SELinux evaluation  & $\sim$ Partial     & Enforcing active (Android~16); boot loop resolved; \\
                    &                    & \quad full policy audit pending \\
FBE~2.0 evaluation  & $\sim$ Partial     & FBE configured and boot loop diagnosed (TC-S-005); \\
                    &                    & \quad end-to-end operation not confirmed within project timeline \\
Reproducible pipeline & \checkmark Pass  & \texttt{rpi5-build} CLI \\
TARA                & \checkmark Pass    & Chapter~\ref{ch:tara}: 23 threats, 10 residual risks \\
\bottomrule
\end{tabular}
\end{table}

% ----------------------------------------------------------------
\section{Limitations}
% ----------------------------------------------------------------

\begin{itemize}
    \item \textbf{FBE post-fix capture}: The \texttt{e2fsdroid} fix that resolves the metadata SELinux boot loop (Section~\ref{sec:impl_fbe_selinux}) is implemented in the build pipeline, however the validation is pending.
    \item \textbf{SELinux policy}: SELinux evaluation revealed that Android~16 transitions to Enforcing mode unconditionally at policy load time, regardless of the \path{androidboot.selinux=permissive} kernel argument. The resulting boot loop was diagnosed and resolved (Section~\ref{sec:impl_fbe_selinux}), confirming basic policy compatibility. However, a full device-specific audit to confirm no residual AVC denials remain in steady state is required before treating Enforcing mode as fully validated \cite{selinux_device_policy}.
    \item \textbf{Hardware-backed keystore}: The RPi5 has no TEE or dedicated security chip. FBE key material is software-protected, which is weaker against a fully privileged attacker with physical access \cite{android_data_encryption}.
    \item \textbf{OTA updates}: The pipeline produces a full SD card image. Incremental OTA support (A/B partitions, \texttt{update\_engine}) is not implemented.
    \item \textbf{Rollback protection}: AVB rollback index counters require tamper-evident storage (eMMC RPMB or TrustZone), neither of which is available on the RPi5 in this configuration. Rollback index values are fixed at zero.
\end{itemize}
